\section{Black White and Grey}

Being a review of some recent, and not so recent, films: \emph{For Greater Glory}, \emph{Swept Away}, \emph{I Am Love}, and \emph{50 Shades of Grey}.

\paragraph{Swept Away}
This 1974 film was directed by \textbf{Lina Wertmuller}, a committed leftist. A small group of Milanese industrialists take a cruise on a private yacht on the Mediterranean, crewed by some Sicilian sailors including the communist Gennarino. Milan is the Nordic part of Italy and the main female passenger, Raffaella, is blonde, lean, and dolichocephalic.

The theme is how a love relationship might work in the state of nature compared one in an the artificial society. The bourgeois passengers are fashionably left, except for Raffaella, who defends her social privilege, often by maltreating or insulting the crew. For example, she chastised Gennarino for overcooking the pasta (that would have irritated me, also).

It becomes clear that the ``progressives" on the cruise have no intention of reversing the social order maintained by financial privilege. Rather, for them, progress consists in overthrowing the established spiritual order of the Church, which was still quite strong in Italy at that time. Hence, they praised the recent plebiscite allowing divorce. Moreover, they looked forward to further erosion of the spiritual authority such as legalized abortion.

One of the passengers gave it away as he gazed upon the beautiful open sea, with no other boats in sight. He said, ``I hope this doesn't get ruined with hordes of poor people." Obviously, he had no inkling of what would someday happen at Lampedusa, the camp of the saints.

While on a side trip away from the yacht, Gennarino and Raffaella get stranded on an uninhabited island. There, in the State of Nature, bourgeois social structures count for nothing. Gennarino is able to survive by capturing lobsters and finding shelter. Raffaella continues to nag him as though nothing had changed, but the elements and hunger finally get to her, so she submits to Gennarino. Through all the time on the island, Raffaella is perfectly coiffed with makeup; that is how she appears to Gennarino, who desires her, but on his terms.

When he captures and skins a rabbit, she initially turns away, calling him ``cruel". Yet, at that same moment she begins to fall in love with him, acquiescing to his demands. He frequently slaps her, but she seems to enjoy it. Now, women have asked me to slap them in the middle of coitus, but it always seemed odd to me, so I don't know if it is a common fantasy.

When she sees a ship come by, she hides in order to remain with Gennarino in that Edenic bower. She confesses to him that her one regret that she could not give up her virginity to him, as if it were the most important gift a woman can give a man. Unexpectedly, she asks to be sodomized instead; apparently her anus was still pure and virginal, not having known man. Although that was the best gift she had to offer him, Gennarino seemed uninterested. Rather, he wanted her undying love, freely given, in all possible worlds.

To test her, Gennarino flagged down the next boat. On leaving the state of nature, Rafaella decides to resume her previous life with her husband and their financial advantages. Gennarino reluctantly goes home with his dowdy and nagging wife.

\paragraph{Degeneration of the Generations}
\emph{I Am Love} (\emph{Io Sono L'Amore}) was a minor hit in 2009. It describes three generations of the Recchi family, who owned a textile manufacture in Milan. The spiritual and moral decline of the family mirrors the degeneration of society:

\begin{enumerate}
\item \textbf{Entrepreneurial}: The founder of the factory and patriarch of the family, Edoardo Sr, was assertive, creative, and self-confident. He describes how he had to manoeuver the business through the Fascist period and thereafter. 
\item \textbf{Managerial}. His son, Tancredi, represents the next phase. He used his resources to collect art. When in Russia, he met Emma, whom he took back to Italy as his wife. He is intelligent and competent to run an established company successfully. However, you never get the sense that he could have founded a company on his own. 
\item \textbf{Decadent}. The grandchildren are spoiled and decadent, incapable of either founding or managing an enterprise. 
\end{enumerate}
Edoardo Sr passes the business on to both Tancredi and Edoardo Jr, one of the grandchildren. Junior has no interest in the enterprise and Tancredi is the real manager. Instead, Junior plans to finance a natural, but seemingly impractical, restaurant with a friend, the chef Antonio.

The sister Elisabetta, a wannabe artist, announces that she has become a lesbian. At the airport on a trip home, she ignores her humiliated boyfriend, left standing with a bouquet of flowers in his hand. When an international conglomerate offers to buy out the family business, she excitedly wonders about how rich she will become after they cash out.

Meanwhile, Antonio, who has a rocky relationship with his own father, starts a weird Oedipal relationship with Emma, the mother of his friend and business partner, Junior. When Junior becomes aware of it, he angrily confronts his mother. During the spat, he slips, hitting his head on the edge of the swimming pool and dies.

After the funeral, Emma confesses her affair to Tancredi, who immediately disowns her. Emma rushes back to the house to pack up some things to leave. Before she goes, her daughter Elisabetta gives her a knowing look of approval.

On the surface, this may seem like a story of liberation. Both Emma and Elisabetta follow their passions, their bliss, into their true happiness. That is what women today call ``chemistry", a passive response to biochemical hormonal surges, that trump any sense of responsibility or prior commitments. Without the question of morality, the third dimension of the story cannot be understood. If everything is ``natural", i.e., two dimensional, then adulterous and lesbian desires are simply natural forces, neither good not evil, whose compulsion need not be resisted.

Rather, these forces are not natural, but are instead telluric, Demetrian, or Aphroditic forces: sexual expression is undifferentiated so the object of one's affections does not matter, whether a man or a woman. Promiscuity is normal. Although Emma's and Elisabetta's privileges are the result of patriarchy, the cuckolded husband and the spurned boyfriend received their just rewards.

\paragraph{Shades of Grey}
\textbf{J J Bachofen} described the reaction of the Hellenic mind when it first learned of earlier matriarchal cultures:

\begin{quotex}
The older system represented an utter puzzle to the patriarchal mind, which consequently could not have conceived any part of it. 

\end{quotex}
I find myself in the same state of puzzlement in trying to understand the phenomenon of \emph{50 Shades of Grey}. With over 50 million copies sold followed by a wildly popular film version, the hetaeric influence reveals itself as stronger than expected. You begin to wonder how many women you may pass by are secret admirers of the book. There was an older woman, who described herself as ``Christian" and ``very conservative", but loved Shades for its ``honesty". I asked her if she would like to engage in that sort of relationship, but she immediately ``unfriended" me without explanation.

Since I was unable to read more than a few pages of the book, I watched the move version instead. Anastasia, a college student studying English literature, goes to interview Christian Grey, a 20 something handsome billionaire, who wants to help the poor people in Africa through various business ventures. The demure Anastasia is overwhelmed by the phallic symbol of Grey's skyscraper. She is flustered during the interview, yet Christian seems intrigued. He explains that he owes his success to his ability to understand people. When he wrongly guesses that Anastasia loved Bronte, rather that Hardy, he must have felt a sense of mystery in Anastasia.

Grey loves motion, owns several cars, even a helicopter and a glider. The bookish Anastasia can't get enough of that motion. Dakota Johnson, as Anastasia, displays her acting talents when in the glider, she captures the same expression of your dog when he rides in the back of your pickup truck.

Grey shows Anastasia his ``playroom", full of various bondage and domination items. Instead of running out the door, she shows some interest. Of course, all acts are to be fully consensual, explicitly mentioned in a binding contract. That seems to be the direction of sexual encounters on college campuses today.

While Grey wants to tie her up, beat her, and lick her body parts, Anastasia wants the full boyfriend experience. Therein lies the whole dramatic tension, if you actually give a damn.

\paragraph{Spiritual Warriors}
\emph{For Greater Glory} is the story of the Cristero War in Mexico in the 1920s. The Masonic government cracks down on the Catholic Church giving rise to an armed rebellion, a revival of the ideal of the spiritual warrior. It is not just Islam that is willing to engage in the lesser Holy War.

The film is better than the reviews would have it. If you don't see it as the story of the creation of saints, amid the horrors of religious persecution, it will mean nothing to you. Stories of saints and martyrs are always edifying.

This is just one particular moment of history, which must include persecutions by communists and the terror of the Vendee. So don't believe it can't happen again in the West. I mean the persecution, because a new generation of spiritual warriors seems unlikely.

One thing, however, is certain. If large bands of Cristeros began to migrate from Mexico to the USA, then liberals would be clamoring to build a wall.

\textbf{Viva Cristo Rey} !

\textbf{Los Cristeros, orate pro nobis}



\flrightit{Posted on 2015-07-16 by Cologero }

\begin{center}* * *\end{center}

\begin{footnotesize}\begin{sffamily}



\texttt{Matt on 2015-07-16 at 15:58 said: }

``Dakota Johnson, as Anastasia, displays her acting talents when in the glider, she captures the same expression of your dog when he rides in the back of your pickup truck."

``While Grey wants to tie her up, beat her, and lick her body parts, Anastasia wants the full boyfriend experience. Therein lies the whole dramatic tension, if you actually give a damn."

Over the years you've provided some great zingers on products and trends of the modern world, and these two rank at the top of the list.


\hfill

\texttt{Sparrow on 2015-07-16 at 21:50 said: }

``Rather, these forces are not natural, but are instead telluric, Demetrian, or Aphroditic forces:"

I'm likely wrong, but wouldn't it be more correct to state that these forces and acts are completely natural (given that they can be observed in the animal world, and the aforemented races of the spirit) while the Traditional man strives to go beyond such base, naturalistic impulses?


\hfill

\texttt{Matt on 2015-07-16 at 22:36 said: }

Sparrow,

I think Cologero was using the term `natural' in the sense of what a thing is (its essence), with the specific context of what a human being is (the human essence) and its essential attributes. I don't think he was using the term in the sense of referring to the wild/natural world, though many that defend the impulses and passions in question will point to the natural world (the wild) as a justification in arguing for them. They would be committing the naturalistic fallacy in that regard.

So in light of Tradition, those passions and inclinations, rather than being natural and arising from the human essence, are in fact unnatural; they are ethical privations (lack of moral goods a human should have given their proper essence) and represent a gap between the human essence and its full/complete manifestation.


\hfill

\texttt{Sparrow on 2015-07-16 at 23:06 said: }

Matt,

Thank you for clarification. I realized that I was likely misinterpreting the use of the word ``natural." but at that time, my message had already been sent.


\hfill

\texttt{obscure on 2015-07-17 at 15:33 said: }

Matt and Sparrow,

In the physiological triad of the soul (vegetative/nutritive, sentient/animal and rational) the passions are indeed natural if only the two lower animating principles are considered. Thomas Aquinas refers to the lower appetites as concupiscence and irascibility. Concupiscence is of course the desire to feed and multiply which man essentially shares with plant-life, but this is enhanced by sentient irascibility which is essentially a substrate of sense-experience which consists of pre-intentional fear and anger; the instinctual desire for safety and preservation of the body; the general arousal relative to sense-objects. The free will is the rational appetite of course and for a rational animal to only operate according to the lower appetities is indeed a degradation of its nature. Before even considering what is above nature (grace) the relative perfection of nature is required on the part of the human subject. The subject must `do what it can' and prepare itself.


\hfill

\texttt{Matt on 2015-07-17 at 19:06 said: }

Obscure,

Thank you for mentioning the concupiscible and irascible appetites and that they should be kept in mind when considering these topics.

I should have been clearer when using the term passions, and so I'll try to provide some needed clarification. By passions, I mean more in the eastern sense of the term, which is `sins'. instead of the scholastic sense referring to our sense-appetites which make up our essence as embodied beings.


\hfill

\texttt{obscure on 2015-07-19 at 01:29 said: }

Passion is simply the opposite of action. Categorically, it is bad. Of course, a finite subject must always be passive in some way and ultimately it is passivity to God alone which is the aim; not through being freed of social obligations, devotion, necessary bodily appetites, etc., but through the `integration' of these things into the one final aim. In other words, one is to play his role perfectly solely for the sake of the Infinite which has granted all things.


\end{sffamily}\end{footnotesize}
